\documentclass[a4paper,12pt]{article}
% https://www.sharelatex.com/learn/Hyperlinks
\usepackage[utf8]{inputenc}
%\usepackage{isolatin1}
\usepackage[french]{babel}
\usepackage{mathtools,amsmath,amssymb,mathrsfs}
\usepackage{hyperref}
\usepackage{amsmath}
\hypersetup{colorlinks=true,
citecolor = blue
} % pour avoir des liens en couleur.
%\usepackage{mathtools,amsmath,amssymb,mathabx,mathrsfs}

\DeclareMathOperator{\sinc}{sinc}

\pagestyle{empty}
\begin{document}

% pour avoir de véritables puces. aplacer apres \begin{document}
\renewcommand{\labelitemi}{\textbullet} 
\newcommand{\tn}{\ensuremath{T_N}}
\newcommand{\ts}{\ensuremath{T_S}}
\newcommand{\tno}{\ensuremath{T_N^O}}
\newcommand{\tso}{\ensuremath{T_S^O}}

\newcommand{\gn}{\ensuremath{g_N}}
\newcommand{\xt}{\ensuremath{x^T}}
\newcommand{\gs}{\ensuremath{g_S}}
\newcommand{\cn}{\ensuremath{C_N}}
\newcommand{\cs}{\ensuremath{C_S}}
\newcommand{\cno}{\ensuremath{C_N^O}}
\newcommand{\cso}{\ensuremath{C_S^O}}
\newcommand{\tauns}{\ensuremath{\tau_{NS}}}
\newcommand{\taunso}{\ensuremath{\tau_{NS}^O}}
\newcommand{\nlat}{\ensuremath{n_{\text{lat}}}}
\section{Définitions}

\subsection{\'{E}tat du système}
L'état du système est défini par:
\begin{itemize}
  \item \tn:  \  température de l'atmosphère et de la couche supérieure de l'océan dans l'hémisphère Nord,
  \item \ts:  \  température de l'atmosphère et de la couche supérieure de l'océan dans l'hémisphère Sud, 
  \item \tno: \  température de la couche profonde de l'océan dans l'hémisphère Nord,                     
  \item \tso: \  température de la couche profonde de l'océan dans l'hémisphère Sud,                      
\end{itemize}


\subsection{Entrées}

\begin{itemize}
  \item $e_{jl}(t)$\ émission de l'acteur $j$ au point d'injection $l$ à l'instant $t$.
  \item $f$:\ forçage radiatif externe (hors aérosols injectés)
\end{itemize}

\subsection{Autres définitions}
On définit par ailleurs les variables suivantes:


\begin{itemize}
  \item $n_a$ le nombre d'acteurs
  \item \nlat le nombre de points d'injections
  \item $e_{l}(t)=\sum_{j} e_{jl}(t)$ émissions cumulées de tous les acteurs au point d'injection $l$ à l'instant $t$.
  \item \cn \ la capacité calorifique de l'atmosphère et de la couche supérieure de l'océan pour l'hémisphère Nord,
  \item \cs \  la capacité calorifique de l'atmosphère et de la couche supérieure de l'océan pour l'hémisphère Sud, 
  \item \cno\  la capacité calorifique de la couche profonde de l'océan pour l'hémisphère Nord,                     
  \item \cso\  la capacité calorifique de la couche profonde de l'océan pour l'hémisphère Sud,                      
  \item \gn \ forçage radiatif dû aux aéosols injectés, pour l'hémisphère Nord ($W$)
  \item \gs \ forçage radiatif dû aux aéosols injectés, pour l'hémisphère Sud ($W$)
%  \item $C_{ON}=0.6$ fraction d'océan pour l'hémisphère Nord
%  \item $f_{OS}=0.8$ fraction d'océan dans l'hémisphère Sud
  \item $h_{S}^{l}(u)e_l(t-u)$  est la contribution à l'épaisseur optique pour l'hémisphère Sud 
         à la date  $t$, des émissions au point $l$  à la date $t-u$.
  \item $h_{N}^{l}(u)e_l(t-u)$  est la contribution à l'épaisseur optique pour l'hémisphère Nord
        à la date $t$,  des émissions au point $l$  à la date $t-u$.
       \item $K_{jl}^p$, vecteur ligne de taille 4, terme proportionnel d'un PID pour l'acteur $j$,
             le point d'émission $l$  
       \item $K_{jl}^d$, vecteur ligne de taille 4, terme dérivé d'un PID pour l'acteur $j$,
             le point d'émission $l$  
       \item $K_{jl}^i$, vecteur ligne de taille 4, terme intégral d'un PID pour l'acteur $j$,
             le point d'émission $l$  
   \item $\Delta t$, temps d'intégration pour le terme intégral des PID
\end{itemize}


\section{Equation régissant l'évolution du système}

Les dérivées des variables définissant l'état du système sont définies par les équations suivantes:

\begin{equation}
  \left\{
\begin{aligned}
  \frac{d\tn}{dt}&=\frac{1}{\cn}\left[f+g_{\text{eff}}g_n - \lambda\tn - \gamma\left(\tn-\tno\right) - \tauns\left(\tn-\ts\right)\right] \\
  \frac{d\tno}{dt}&=\frac{1}{\cso}\left[-\gamma\left(\tno-\tn\right)-\tauns\left(\tno-\tso\right)\right] \\
  \frac{d\ts}{dt}&=\frac{1}{\cs}\left[f+g_{\text{eff}}g_s - \lambda T_s - \gamma\left(\ts-\tso\right) - \tauns\left(\ts-\tn\right)\right] \\
  \frac{d\tso}{dt}&=\frac{1}{\cso}\left[-\gamma\left(\tso - \ts\right) - \tauns\left(\tso-\tno\right)\right] \\
\end{aligned}
\right.
\end{equation}

où

\begin{align*}
  g_S(t)&=A_S\int h_S(u)\tau_S(t-u) du\\
     &=A_S\int h_S^l(u)\sum_{j=1}^{n_a}\sum_{l=1}^{n_{\text{lat}}}  e_{jl}(t-u) du \\
     &=A_S\sum_{j=1}^{n_a}\sum_{l=1}^{n_{\text{lat}}} \int h_S^l(u) e_{jl}(t-u) du \\
     &=A_S\sum_{j=1}^{n_a}\sum_{l=1}^{n_{\text{lat}}}  \left(h_S^l \ast  e_{jl}\right)(t)
\end{align*}

 et 
\begin{equation}
  g_N(t) =A_N\sum_{j=1}^{n_a}\sum_{l=1}^{n_{\text{lat}}}  \left(h_N^l \ast  e_{jl}\right)(t)
\end{equation}

en changeant l'ordre des termes et en développant les termes dûs au forçage des aérosols, ce système devient:
\begin{equation}
  \left\{
  \begin{aligned}
  \frac{d\tn}{dt}&=  -\frac{\lambda+\gamma+\tauns}{\cn}\tn +\frac{\gamma}{\cn}\tno + \frac{\tauns}{\cn}\ts  
                  +A_N\sum_{j=1}^{n_a}\sum_{l=1}^{n_{\text{lat}}}  \left(h_N^l \ast  e_{jl}\right)(t)
                  +\frac{f}{C_N} \\
  \frac{d\tno}{dt}&=\frac{\gamma}{\cno}\tn -\frac{\gamma+\tauns}{\cno}\tno +\frac{\tauns}{\cno}\tso \\
  \frac{d\ts}{dt}&=-\frac{\lambda+\gamma+\tauns}{\cs}\ts +\frac{\gamma}{\cs}\tso
                  +\frac{\tauns}{\cs}\tn 
                  +A_S\sum_{j=1}^{n_a}\sum_{l=1}^{n_{\text{lat}}}  \left(h_S^l \ast  e_{jl}\right)(t)
                  +\frac{f}{C_S} \\
  \frac{d\tso}{dt}&=\frac{\gamma}{\cso}\ts -\frac{\gamma+\tauns}{\cso}\tso +\frac{\tauns}{\cso}\tno \\
\end{aligned}
\right.
\end{equation}

on peut le réécrire sous une forme condensée:

\begin{equation} \dot x = Ax + B  + Df\end{equation}

  avec 
\[
  A = \begin{bmatrix}
    -\frac{\lambda+\gamma+\tauns}{\cn} & \frac{\gamma}{\cn} &   \frac{\tauns}{\cn} & 0  \\
    \frac{\gamma}{\cno} &  -\frac{\gamma+\tauns}{\cno} & 0 & \frac{\tauns}{\cno} \\
    \frac{\tauns}{\cs} & 0 & -\frac{\lambda+\gamma+\tauns}{\cs} & \frac{\gamma}{\cs} \\
    0 & \frac{\tauns}{\cso} &     \frac{\gamma}{\cso} & -\frac{\gamma+\tauns}{\cso}   \\
\end{bmatrix}\]

\[
  x =
  \begin{bmatrix}
  \tn \\
  \tno \\
  \ts \\
  \tso 
\end{bmatrix}\]

\[
  D =
  \begin{bmatrix}
  1/\cn \\
  0 \\
  1/\cs \\
  0 \\
\end{bmatrix}\]



\[
  B =  \sum_{j=1}^{n_a}\sum_{l=1}^{n_{\text{lat}}} 
  \begin{bmatrix}
 A_N \left(h_N^l \ast  e_{jl}\right)(t)  \\
0   \\
 A_S \left(h_S^l \ast  e_{jl}\right)(t)  \\
0  \\
\end{bmatrix}\]


\section{Correction}

Soit un système dont l'état est décrit par un vecteur $x$, qui est déterminé via un modèle par un vecteur d'entrée $e$. On cherche
à ajuster les entrées pour obtenir un état $x_c$. 


\subsection{Correcteur PID simple}

Ce correcteur détermine la valeur  d'une seule composante $e_j$ du vecteur d'entrée à partir d'une seule composante $x_l$ du vecteur d'état, à l'instant $t$, 
de ses valeurs dans le passé\ $x_l(t')\ t' < t$, et de sa valeur cible $x^c_{l}$

\begin{equation} 
  e_{j}=-K^p(x_l-x_l^c) - K^d\frac{d(x_l-x^c_l)}{dt} - K^i\int_0^{t} \left(x_l(t')-x^c_l\right)dt'
\end{equation}

\subsection{Correcteur PID multiple}

Ce correcteur détermine le vecteur $e$ à partir du vecteur d'état $x$ à l'instant $t$, 
de ses valeurs dans le passé\ $x(t')\ t' < t$, et de sa valeur cible $x^c$. Contrairement au cas précédent,
on ce correcteur peut utiliser toutes les composantes du vecteur d'état. 

Notons $n_s$ la taille du vecteur d'état, et $n_c$ la taille du vecteur d'entrée.

Pour tout $l$, $1 \le j \le n_s$:
\begin{equation} 
  e_{j}=-\sum_{l=1}^{n_s}\alpha_l \left[K_{jl}^p(x_m-x_l^c) + K_{jl}^d\frac{d(x_l-x^c_l)}{dt} + K^i_{jl}\int_0^{t} (x_l(t')-x^c_l)dt' \right]
\end{equation}

Les $\alpha_k$ permettent de donner des poids différents aux différentes composantes du vecteur d'état.


\subsection{Correcteur PID multiple, avec fonction de poids pour le terme intégral}

Afin de faire varier le poids des différents temps passés dans le terme intégral $K_i \ldots$, on peut utiliser une fonction de poids $\alpha$:

\begin{equation} 
  e_{j}(t)=-\sum_{l=1}^{n_s}\alpha_l \left[K_{jl}^p(x_m-x_l^c) + 
                                           K_{jl}^d\frac{d(x_l-x^c_l)}{dt} +
                                           K^i_{jl}\int_0^{t}\alpha(u) (x_l(u)-x^c_l(u))du \right]
\end{equation}

Cela permet 
\begin{itemize}
  \item soit de limiter la mémoire de la variable $x$ ($\alpha(u)=0$\ pour\  $u\le t_1$ ou $u\le t-\Delta u$
  \item soit de considérer un retard à l'action ($\alpha(u)=0$\ si\ $u\ge t-\Delta u$
\end{itemize}
\end{document}
